% TOP_PROBLEM: {"id": 2857, "title": "Kirby Problem 3.59", "category_id": 7, "category": {"id": 7, "name": "topology", "display_name": "Topology", "description": "Properties preserved under continuous deformations.", "slug": "topology", "order_index": 7, "created_at": "2026-07-31T15:26:25.671Z"}, "status": "open"}
% FIRST_POSTED: null
% ENTRY_KIND: statement_only
% Imported from: batch12.tex; UnsolvedMath ID 2857
\documentclass[11pt]{article}
\usepackage[margin=25mm]{geometry}
\usepackage{fontspec}
\setmainfont{Latin Modern Roman}
\usepackage{amsmath,amssymb,mathtools}
\usepackage{unicode-math}
\setmathfont{Latin Modern Math}
\usepackage{xurl,hyperref}
\hypersetup{colorlinks=true,urlcolor=blue}
\setcounter{secnumdepth}{0}
\setlength{\parindent}{0pt}
\setlength{\parskip}{5pt}
\setlength{\emergencystretch}{3em}
\newcommand{\sref}[2]{\hyperlink{source:#1:#2}{[S#2]}}
\newcommand{\cataloguescope}[1]{\paragraph{Recorded target.}#1}
\begin{document}
% UnsolvedMath ID 2857; source catalogue code KP-3.59
\setcounter{section}{2856}
\section{Kirby Problem 3.59}\label{2857}
{\small\sffamily UnsolvedMath ID 2857 · Topology}
\subsection{Definitions and mathematical statement}

\paragraph{Shared notation.}$\mathbb N=\{1,2,\ldots\}$, and $\mathbb Z,\mathbb Q,\mathbb R,\mathbb C$ denote the integers, rationals, reals and complex numbers. Any other conventions are specified locally.


A balanced sutured manifold $(M,\gamma)$ is a compact oriented $3$-manifold without closed components, with oriented sutures in annular neighborhoods $\gamma\subset\partial M$, whose complement is $R_+(\gamma)\sqcup R_-(\gamma)$, with compatible boundary orientations, $\chi(R_+)=\chi(R_-)$, and every boundary component meeting $\gamma$. The group $\operatorname{SFH}(M,\gamma;k)$ is sutured Heegaard Floer homology, the homology of its admissible sutured Heegaard chain complex. The group $\operatorname{SHI}(M,\gamma;\mathbb C)$ is sutured instanton homology: the distinguished surface-eigenspace of instanton Floer homology of an admissible closure of the sutured manifold. Their isomorphism classes are independent of diagram and closure choices. For a closed oriented $Y$, $I^\#(Y;\mathbb C)$ is framed instanton Floer homology and $\widehat{\operatorname{HF}}(Y;\mathbb C)$ is hat Heegaard Floer homology. The corresponding knot versions arise from the knot exterior with meridional sutures.
\paragraph{Statement recorded in the supplied source.}
For every balanced sutured manifold $(M,\gamma)$, prove an isomorphism of complex vector spaces
\[
\operatorname{SHI}(M,\gamma;\mathbb C)\cong\operatorname{SFH}(M,\gamma;\mathbb C).
\]
This includes the resulting closed-manifold comparison $I^\#(Y;\mathbb C)\cong\widehat{\operatorname{HF}}(Y;\mathbb C)$ and the corresponding non-equivariant knot-invariant comparison. This is the explicit sutured interpretation given in the source. \sref{2857}{2}
\subsection{Short English statement}
Prove that sutured instanton and sutured Heegaard Floer homology agree over the complex numbers for every balanced sutured manifold, including the closed-manifold and knot special cases.
\subsection{Sources}
\begin{description}
\item[\hypertarget{source:2857:1}{S1}] ULAM.AI and the UnsolvedMath Contributors, \emph{UnsolvedMath: A Curated Collection of Open Mathematics Problems}, record 2857 (KP-3.59). CC BY 4.0. \url{https://huggingface.co/datasets/ulamai/UnsolvedMath}.
\item[\hypertarget{source:2857:2}{S2}] R. İnanç Baykur, Robion C. Kirby and Daniel Ruberman (editors), \emph{K3: A New Problem List in Low-Dimensional Topology}, Mathematical Surveys and Monographs 295 (2026), Problem 3.59, p. 173 and following remarks; original author version. \url{https://math.berkeley.edu/sites/default/files/surv-295-ruberman-watermarked-author-pdf.pdf}.
\end{description}
\label{end:2857}
\end{document}
